Let \(v\in K_{n,l}\). By the definition of \(A_{n,l}\), for every \(\tau=(\delta,\eta)\in\mathcal T_{n,l}\),
\[
0=(A_{n,l}^{\mathsf T}v)_\tau=\sum_{z=1}^l v_{(\eta(\delta(z)),\delta(z),z)}.
\]
Fix \(z\in\mathcal Z_l\) and \(y_0,y_1\in\mathcal Y_n\), and choose \(\eta(0)=y_0\) and \(\eta(1)=y_1\). Taking \(\delta(j)=0\) for every \(j\) gives
\[
\sum_{j=1}^l v_{(y_0,0,j)}=0.
\]
Taking instead \(\delta(z)=1\) and \(\delta(j)=0\) for \(j\ne z\) gives
\[
v_{(y_1,1,z)}+\sum_{j\ne z}v_{(y_0,0,j)}=0.
\]
Subtracting yields \(v_{(y_1,1,z)}=v_{(y_0,0,z)}\). Since \(y_0,y_1\) were arbitrary, all \(2n\) coordinates in arm \(z\) have a common value \(a_z\). The all-zero treatment word then gives \(\sum_z a_z=0\). Conversely, if \(v_{(y,d,z)}=a_z\) and \(\sum_z a_z=0\), then
\[
(A_{n,l}^{\mathsf T}v)_\tau=\sum_{z=1}^l a_z=0
\]
for every type \(\tau\). Thus the displayed description of \(K_{n,l}\) holds. The map from \(\{a\in\mathbb R^l:\sum_z a_z=0\}\) to \(K_{n,l}\) is a linear bijection, so \(\dim K_{n,l}=l-1\).

For \(p\in\mathcal C_{n,l}\), armwise normalization gives
\[
v^{\mathsf T}p=\sum_{z=1}^l a_z\sum_{y\in\mathcal Y_n}\sum_{d\in\mathcal D}p_{y,d,z}=\sum_{z=1}^l a_z=0.
\]
For \(\lambda\in\mathcal H_{n,l}\), put \(g(\lambda)=\Gamma(\lambda)-\lambda\). Its arm constants are
\[
g_z(\lambda)=-\lambda_{(0,0,z)}\quad(z\ge2),\qquad g_1(\lambda)=\sum_{z=2}^l\lambda_{(0,0,z)},
\]
so \(g(\lambda)\in K_{n,l}\). Hence \(A_{n,l}^{\mathsf T}\Gamma(\lambda)=A_{n,l}^{\mathsf T}\lambda\), and \(\Gamma(\lambda)\in\mathcal H^0_{n,l}\). If \(\lambda'=\lambda+k\) with \(k\in K_{n,l}\) having arm constants \(a_z\), then \(g_z(\lambda')=g_z(\lambda)-a_z\) for \(z\ge2\), while
\[
g_1(\lambda')=g_1(\lambda)+\sum_{z=2}^l a_z=g_1(\lambda)-a_1.
\]
Thus \(\Gamma(\lambda')=\Gamma(\lambda)\). Conversely, if \(\lambda+k\in\mathcal G_{n,l}\) with \(k\in K_{n,l}\), its constants must satisfy \(a_z=-\lambda_{(0,0,z)}\) for \(z\ge2\) and \(a_1=-\sum_{z=2}^l a_z\), so \(\lambda+k=\Gamma(\lambda)\). Every equivalence class therefore meets \(\mathcal G_{n,l}\) exactly once.

The lineality space of the section is
\[
\operatorname{lin}(\mathcal H^0_{n,l})=K_{n,l}\cap\mathcal G_{n,l}.
\]
An element of this intersection has \(a_z=0\) for \(z\ge2\) by the gauge equations, and then \(a_1=0\) because \(\sum_z a_z=0\). Hence the intersection is \(\{0\}\), proving that \(\mathcal H^0_{n,l}\) is pointed. Finally, \(\Gamma\) induces an affine bijection from the quotient polyhedron \(\mathcal H_{n,l}/K_{n,l}\) onto \(\mathcal H^0_{n,l}\). An affine bijection preserves extreme points, giving the claimed bijection between section vertices and M-distinct quotient vertices.